搜索吞吐采用\cref{fig:3-boundary} 中同步的主机到主机边界，计入查询暂存、传输和预处理，例如 JHQ 的旋转与 CAGRA int8 的查询量化。各方法均通过页锁定（pinned）主机缓冲上传查询；JHQ 将结果返回普通可换页（pageable）内存，基线返回页锁定内存。JHQ／RaBitQ 由 C++ 调用，CAGRA／IVF-PQ 由 Python 调用，句柄和输出缓冲复用。这些接口成本计入测量，对小 batch 和接近持平的比较尤其重要。

\par\medskip\noindent\begin{minipage}{\linewidth}
\centering\captionsetup{hypcap=false}
\resizebox{\linewidth}{!}{% Method-neutral search-call boundary; not a stage-time plot.
\begin{tikzpicture}[
 font=\small\linespread{1}\selectfont,
 box/.style={draw=ThesisNavy!65,rounded corners=2pt,align=center,
 text width=2.9cm,minimum height=1.7cm,inner sep=4pt},
 arr/.style={-{Latex[length=2mm]},draw=ThesisNavy!80,thick}]
\node[box] (host) at (0,0) {\mbox{Host FP32 queries}\\\mbox{Pinned staging}};
\node[box,fill=ThesisNavy!7] (upload) at (3.65,0)
 {\mbox{H2D transfer}\\Query\\\mbox{preprocessing}};
\node[box,fill=ThesisNavy!7] (gpu) at (7.3,0)
 {GPU search\\kernels\\\mbox{(method-specific)}};
\node[box] (result) at (10.95,0)
 {D2H results\\GPU\\\mbox{synchronisation}};
\draw[arr] (host) -- (upload);
\draw[arr] (upload) -- (gpu);
\draw[arr] (gpu) -- (result);
\draw[ThesisNavy,thick] (-1.6,1.15) -- (-1.6,1.45)
 -- (12.55,1.45) -- (12.55,1.15);
\node[fill=white,inner sep=3pt] at (5.475,1.45)
 {Search timing: host to host};
\node[box,text width=6.2cm,minimum height=1.5cm] (setup) at (1.75,-2.45)
 {Outside search timing\\Data loading; model training/loading\\Index construction; warm-up};
\node[box,text width=6.2cm,minimum height=1.5cm] (cal) at (9.2,-2.45)
 {Budget calibration (Chapter 6)\\Timed separately from search};
\draw[densely dashed,ThesisNavy!60] (-1.6,-1.1) -- (12.55,-1.1);
\end{tikzpicture}
}
\captionof{figure}[主机与设备计时边界]{计时从主机内存中的 FP32 查询开始，到结果返回主机结束，包含输入暂存、主机到设备（host-to-device，H2D）传输、查询预处理、GPU 搜索、设备到主机（device-to-host，D2H）传输及同步。构建、预热和预算校准不计入；构建与校准成本另行报告。GPU 各阶段时间来自单独的剖析运行，其和不一定等于这一端到端时间。}
\label{fig:3-boundary}
\end{minipage}\par\medskip

主比较 使用大小为 1,000 的 batch，预热后测量五轮，通常每轮 24 次调用。对大小为 $b$ 的 batch，第 $r$ 轮有 $m_r$ 次调用，每次耗时 $t_{ri}$ 秒。先用 batch 大小除以轮内平均调用时间，再取各轮吞吐的中位数：
\begin{equation}
 \bar t_r=\frac1{m_r}\sum_{i=1}^{m_r}t_{ri},
 \qquad Q_r=\frac b{\bar t_r},
 \qquad Q=\operatorname{median}_r Q_r.
 \label{eq:3-throughput}
\end{equation}
均值保留有效的慢调用。结果衡量同步 batch 吞吐，而非并发请求的尾延迟。各轮共享同一索引；独立冷构建用于评估训练波动（\cref{sec:7-7}）。预热与运行准备细则见 \cref{sec:a-e1-protocol}。

主比较 使用 RTX 5090（32,607\,MiB，功率上限 575\,W）、驱动 595.71.05、CUDA 13.0 和 cuVS 26.08.01。OpenMP 使用 16 个线程，环境提供 25 个逻辑 CPU。JHQ 内核以 compute\_80 PTX 为目标，由驱动在计时前编译为该 GPU 的指令 \cite{cuda_platform}。GPU 时钟与主机后台活动未受控制，可能带来计时波动。补充实验在后续会话中运行；各项报告说明哪些方法一同计时、是否复用训练模型，以及采用何种数值精度（\cref{sec:a-supplement,sec:4-2,sec:7-precision}）。

显存统计区分索引分配、包含上下文和工作区的搜索时进程总占用，以及构建峰值。JHQ 的净值扣除背景分配，CAGRA 的数值则衡量索引分配，二者边界不同（\cref{sec:4-4,sec:7-7}）。构建峰值每 20\,ms 采样，可能遗漏更短的瞬时峰值；构建时间覆盖 train／add 调用，不含输入载入。
